\subsection{What the Familiar Cartesian Forms Let Us Read}

\begin{corebox}[Choose the representation adapted to the question]
The parametric, normal, and slope forms are not competing definitions. They
expose different information about the same geometric line, so the useful form
depends on what must be calculated or interpreted.
\end{corebox}

\subsubsection*{Intersections with the coordinate axes}

\begin{derivationbox}[Impose both defining conditions]
For a nonvertical line
\[
y=mx+b,
\]
setting $x=0$ imposes the condition for the $y$-axis and gives $(0,b)$. If
$m\neq0$, setting $y=0$ imposes the condition for the $x$-axis and gives
\[
x=-\frac bm.
\]
\end{derivationbox}

\begin{theorembox}
For $y=mx+b$,
\[
\boxed{L\cap\{x=0\}=\{(0,b)\}},
\]
and, when $m\neq0$,
\[
\boxed{L\cap\{y=0\}=\left\{\left(-\frac bm,0\right)\right\}}.
\]
\end{theorembox}

\begin{interpretationbox}[General intersection principle]
To find an intersection of two sets, impose both defining conditions at once.
Axis intercepts are the simplest example of this general method.
\end{interpretationbox}

\begin{examplebox}
For
\[
y=-\frac12x+3,
\]
the $y$-axis intersection is $(0,3)$. Setting $y=0$ gives $x=6$, so the
$x$-axis intersection is $(6,0)$.
\end{examplebox}

\subsubsection*{The line through two points}

\begin{derivationbox}[Extracting a direction from two points]
Let
\[
P=(x_1,y_1),\qquad Q=(x_2,y_2),\qquad P\neq Q.
\]
The displacement
\[
Q-P=(x_2-x_1,y_2-y_1)
\]
is nonzero and therefore provides a direction vector of the unique line through
the points.
\end{derivationbox}

\begin{theorembox}
The unique line through distinct points $P$ and $Q$ is
\[
\boxed{L=\{P+t(Q-P):t\in\mathbb R\}}.
\]
This formula handles vertical and nonvertical lines uniformly.
\end{theorembox}

\begin{derivationbox}[Recovering slope notation]
If $x_1\neq x_2$, the horizontal component of $Q-P$ is nonzero, so
\[
\boxed{m=\frac{y_2-y_1}{x_2-x_1}}
\]
and
\[
\boxed{y-y_1=m(x-x_1)}.
\]
If $x_1=x_2$, the line is vertical and is described by
\[
\boxed{x=x_1}.
\]
\end{derivationbox}

\begin{remarkbox}
The exceptional vertical case belongs only to slope notation. The geometric line
and its parametric description require no exception.
\end{remarkbox}

\subsubsection*{Parallel lines in slope notation}

\begin{proofbox}[Why equal slopes mean equal directions]
The nonvertical lines
\[
y=m_1x+b_1,
\qquad
y=m_2x+b_2
\]
have direction vectors $(1,m_1)$ and $(1,m_2)$. These directions are parallel
exactly when $m_1=m_2$.
\end{proofbox}

\begin{theorembox}
For distinct nonvertical lines,
\[
\boxed{
L_1\parallel L_2
\iff
m_1=m_2\text{ and }b_1\neq b_2
}.
\]
If also $b_1=b_2$, the two equations describe the same line.
\end{theorembox}

\begin{remarkbox}
Slope formulas are efficient coordinate consequences of the vector
construction. They should be used when they simplify a question, but not at the
cost of hiding the more general descriptions that handle vertical lines and
extend directly to space.
\end{remarkbox}
